% Manual for the teudot package. Compile with XeLaTeX (twice).
\documentclass[11pt]{article}
\usepackage{shiur}
\usepackage{teudot}
\usepackage{fancyvrb}
\usepackage{booktabs}
\setshiurfont{hebrew}{FreeSerif}
\setshiurfont{mikra}{FreeSerif}
\hypersetup{hidelinks}
\newfontfamily{\codefont}{FreeMono}
\newfontfamily{\hebrewfonttt}{FreeMono}[Script=Hebrew]
\fvset{formatcom=\codefont\small, frame=single, framesep=2mm}
\newcommand{\cs}[1]{\texttt{\char`\\#1}}
\setlength{\emergencystretch}{2em}

\declareteudah[name=Elohist]{e}{red!80!black}
\declareteudah[name=Yahwist]{j}{blue}
\declareteudah[name=Priestly]{p}{green!50!black}

\title{The \texttt{teudot} package\\[0.3em]\large Colouring and isolating the sources of a text}
\author{Gabriel Gendler Yom-Tov\\\small\texttt{shiurtex@thm.mozmail.com}}
\date{Version 0.1, 30 October 2026}

\begin{document}

\maketitle

\section{Introduction}

\texttt{teudot} (\texthebrew{תעודות}, ``documents'', as in \texthebrew{תורת התעודות}, the
documentary hypothesis) is for teaching biblical criticism. A text is marked once with the
sources, or strands, it is divided into, and each copy of it then chooses what to show: every
strand in its own colour, one or two strands alone, or one strand with the others invisible but
still taking their space, so that the versions of a text can be laid over one another.

The examples in this manual use a simplified English passage from Exodus 14 and an illustrative
division; the division of a real text is the author's.

\begin{mekor}[\tanach{exodus}{14:19-23}]
\begin{mekorlang}{english}{Every strand in its colour}
\pasuk[14:19] \teudah{e} The angel of God moved. \teudah{j} The pillar of cloud moved. \pasuk \teudah{e} It came between the camps. \teudah{j} There was cloud and darkness. \pasuk \teudah{p} Moses stretched out his hand. \teudah{j} An east wind blew all night. \teudah{p} The waters were split. \pasuk Israel went into the sea. \pasuk \teudah{j} The Egyptians pursued them.
\end{mekorlang}
\end{mekor}
\begin{mekor*}
\teudotsetup{show=j, others=omit}
\begin{mekorlang}{english}{J alone}
\pasuk[14:19] \teudah{e} The angel of God moved. \teudah{j} The pillar of cloud moved. \pasuk \teudah{e} It came between the camps. \teudah{j} There was cloud and darkness. \pasuk \teudah{p} Moses stretched out his hand. \teudah{j} An east wind blew all night. \teudah{p} The waters were split. \pasuk Israel went into the sea. \pasuk \teudah{j} The Egyptians pursued them.
\end{mekorlang}
\end{mekor*}
\begin{mekor*}
\teudotsetup{show=j, others=ghost}
\begin{mekorlang}{english}{J in its place, the rest invisible}
\pasuk[14:19] \teudah{e} The angel of God moved. \teudah{j} The pillar of cloud moved. \pasuk \teudah{e} It came between the camps. \teudah{j} There was cloud and darkness. \pasuk \teudah{p} Moses stretched out his hand. \teudah{j} An east wind blew all night. \teudah{p} The waters were split. \pasuk Israel went into the sea. \pasuk \teudah{j} The Egyptians pursued them.
\end{mekorlang}
\end{mekor*}

The package works in any XeLaTeX or LuaLaTeX document. With \texttt{shiur}, it marks the text
of a \texttt{mekorlang} directly, and keeps verse numbers in place in every version (see
``\nameref{sec:shiur}'').

\section{Declaring strands}

\begin{Verbatim}
\usepackage{teudot}
\declareteudah[name=Elohist]{e}{red!80!black}
\declareteudah[name=Yahwist]{j}{blue}
\declareteudah[name=Priestly]{p}{green!50!black}
\end{Verbatim}
\cs{declareteudah}\verb|[name=...]{key}{colour}| declares a strand. The key is any short word,
in any case; the colour is any colour \texttt{xcolor} understands. For named colours such as
\texttt{Mahogany}, load \texttt{xcolor} with \texttt{dvipsnames} before \texttt{teudot}. A later
declaration of the same key replaces the earlier one.

\cs{teudotlegend} prints each strand's key in its colour, with its name:
\begin{center}
\teudotlegend
\end{center}

\section{Marking a text}

\cs{teudah}\verb|{key}| switches to a strand until the next \cs{teudah}, like \cs{color}. Text
before the first marker belongs to no strand and always shows, in the colour around it.
Outside \texttt{shiur}, the marked text goes in a \texttt{teudot} environment:
\begin{Verbatim}
\begin{teudot}[show=e, others=omit]
\teudah{e} The angel of God moved. \teudah{j} The pillar of cloud moved.
\end{teudot}
\end{Verbatim}
The package finds the markers by reading the text, so they must be written out in it, at its
top level: not inside a macro, and not inside braces such as \verb|\textbf{...}|. Put them
between words. Inside a pointed Hebrew word, a change of colour can separate a letter from its
vowels.

\section{Choosing a version}

The environment's options, or \cs{teudotsetup}\verb|{...}| for everything after it in the
current group, say what to show:
\begin{description}
\item[\texttt{show=\{e,j\}}] the strands shown in their colours; all of them, if not given.
\item[\texttt{others=}] what happens to the rest:
	\begin{description}
	\item[\texttt{show}] shown in their colours too (the default);
	\item[\texttt{grey}] in light grey, to pick out one strand against the rest;
	\item[\texttt{ghost}] invisible, but taking their space;
	\item[\texttt{omit}] left out, with an ellipsis where text was.
	\end{description}
\item[\texttt{colours=false}] sets the strands shown in the colour around them, for a plain
	text of one strand.
\item[\texttt{greycolour=}, \texttt{ghostcolour=}] the colours for \texttt{grey} (default
	\texttt{black!30}) and \texttt{ghost} (default the page's: white, or on slides the
	background's).
\end{description}
Each copy of a text is set the way its own options say; nothing produces versions by itself. To
show a passage in colour and then each strand alone, include it once for each.

\subsection{Overlays}

With \texttt{others=ghost}, every version of a text has the same lines and the same layout, so
an E version and a J version printed on transparencies, or shown one after the other on slides,
lie exactly over each other. Hidden text is set in the page's colour rather than left out, which
is what keeps the lines the same; it is still in the PDF, and copying text from the page copies
it too.

\section{With \texttt{shiur}}\label{sec:shiur}

In a \texttt{shiur} document the markers go straight into a \texttt{mekorlang}'s text, and
\cs{teudotsetup} inside the \texttt{mekor} chooses its version:
\begin{Verbatim}
\begin{mekor}[\tanach{exodus}{14:19-23}]
\teudotsetup{show=j, others=omit}
\begin{mekorlang}{english}
\pasuk[14:19] \teudah{e} The angel of God moved. \teudah{j} The pillar ...
\end{mekorlang}
\end{mekor}
\end{Verbatim}
A translation takes its own markers. On slides, \cs{mekorbreak} works as usual, and a strand
carries on across it.

Verse numbers belong to no strand: they stay in the colour of the text around them in every
version, coloured, grey or ghosted. So do ellipses, whether teudot's own or a \cs{dilug} written
in a strand's text: a \cs{dilug} stays visible wherever its strand's text is set, even when that text is ghosted,
and goes with it when the strand is left out. When text is left out, a verse shows its number only if
some of the verse shows, and an ellipsis marks each gap, after the number when the start of the
verse is left out (``(20) \ldots{} There was''), and before it when whole verses, or the end of
the verse before, are (``\ldots{} (23) The Egyptians''). A text that begins in the middle of
what it shows starts at the first verse it shows. The J text above comes out as
\begin{quote}
(14:19) ... The pillar of cloud moved. (20) ... There was cloud and darkness. (21) ... An east
wind blew all night. ... (23) The Egyptians pursued them.
\end{quote}
and the P text as
\begin{quote}
(21) Moses stretched out his hand. ... The waters were split. (22) Israel went into the sea.
...
\end{quote}

\section{For package writers}

A package that typesets marked text can use:
\begin{description}
\item[\cs{teudot@part}\texttt{\{true\}\{text\}}] typesets a text with its markers, or a part of
	it, with \texttt{false} for every part but the last (\texttt{shiur} uses it for the parts of a
	source split at \cs{mekorbreak}). \cs{teudot@reset} starts a new text.
\item[\cs{teudot@setmarkers}\texttt{\{regex\}}] says what a verse marker looks like, as an
	\texttt{l3regex} (\texttt{shiur} gives \cs{pasuk} and \cs{perek}); markers then belong to no
	strand.
\item[\cs{teudot@ellipsis}] marks text left out (\texttt{shiur} makes it \cs{dilug}).
\item[\cs{teudot@plain}\texttt{\{text\}}] sets text in the colour around the strands.
\item[\cs{ifteudot@silent}] is true while text left out runs, in a box that is thrown away so
	that counters still advance.
\end{description}

\section{Licence}

Copyright 2026 Gabriel Gendler Yom-Tov. This work may be distributed and/or modified under the
conditions of the \LaTeX{} Project Public License, either version 1.3c of this license or (at
your option) any later version. The latest version of this license is in
\url{https://www.latex-project.org/lppl.txt}.

\end{document}
